\subsection{An explicit gap at every prefix length}
The central additional step is a quantitative bound on the singular defect, with no extrapolation from small matrices.
\begin{theorem}[All-prefix crossing bound]
For $u=\cot\theta>0$ and $n\ge1$, define
\[
 B_n=\frac{u^2+4n^2}{u},\qquad K_n=\frac{B_n^n[(u+2n)^n+1]}{n!},\qquad
 g_n=\frac6{n(n+1)(2n+1)K_n^2}.
\]
Then
\begin{equation}
 \gamma_n:=1-\norm{T_n|_{\mathbf1^\perp}}_{2\to2}^2\ge g_n>0.\label{eq:allgap}
\end{equation}
The identity is the only zero mode of $I-T_nT_n^*$. At fixed $\theta$, $\log g_n=-4n\log n+O_\theta(n)$.
\end{theorem}
\begin{proof}
Use normalized Hilbert--Schmidt norms unless an operator norm is specified. Define
\[
 M_n(z)=(zI+iP_{S,n})\cdots(zI+iP_{S,1}),\qquad U^{(n)}=(u^2+1)^{-n/2}M_n(u).
\]
For $\mathcal C(A)=(\sigma_y\otimes I)A^{t_S}(\sigma_y\otimes I)$, the local identity $\mathcal C(P)=I-P$ implies
\begin{equation}
 \mathcal C[M_n(z-i)]=(z^2+1)^nM_n(z)^{-1}.\label{eq:cross}
\end{equation}
The bath entries in distinct local factors commute, justifying their reversed order under source transpose. No general multiplicativity of partial transpose is assumed. The map $\mathcal C$ preserves Hilbert--Schmidt norm.

Take a traceless bath operator $X$ of norm one, set $Y=T_n(X)$, and write $\widehat X=I_S\otimes X$, $\widehat Y=I_S\otimes Y$. Source conditional expectation is an orthogonal projection, so
\[
 \delta^2:=\norm{U^{(n)}\widehat XU^{(n)\dagger}-\widehat Y}^2=1-\norm{T_n(X)}^2.
\]
For $E(z)=\widehat YM_n(z)-M_n(z)\widehat X$, one has $\norm{E(u)}=(u^2+1)^{n/2}\delta$. If $\widetilde E$ interchanges $X,Y$, crossing gives
\[
 \widetilde E(z-i)=-(z^2+1)^n\mathcal C[M_n(z)^{-1}E(z)M_n(z)^{-1}].
\]
Thus the pair alternates under each shift, returning after two shifts. This matters because norm preservation initially gives an intertwiner, not a commutant.

At $z=u-ik$, each local inverse has operator norm at most $[u^2+(k-1)^2]^{-1/2}$. The norm multiplier per shift is at most
\[
 \left[\frac{u^2+(k+1)^2}{u^2+(k-1)^2}\right]^{n/2}.
\]
These factors telescope. At the $n+1$ even-shift points, $\norm{E(u-2ij)}\le B_n^n\delta$, $j=0,\ldots,n$.

Interpolate the degree-$n$ polynomial $E$. The sum of absolute leading-coefficient weights is $\sum_{j=0}^n[2^nj!(n-j)!]^{-1}=1/n!$. The leading coefficient is $\widehat Y-\widehat X$; the interpolation weights at zero have numerators bounded by $(u+2n)^n$. Hence
\[
 \norm{Y-X}\le B_n^n\delta/n!,\qquad
 \norm{E(0)}\le B_n^n(u+2n)^n\delta/n!.
\]
Since $M_n(0)$ is a unit phase times the cyclic permutation $C$ of all $n+1$ sites, $\norm{[\widehat X,C]}\le K_n\delta$.

Let $P_j$ be the normalized partial-trace conditional expectation making site $j$ identity, where $j=0$ is the source. These projections commute. Since $P_0\widehat X=\widehat X$, cyclic rotation and telescoping give $\norm{(I-P_j)\widehat X}\le jK_n\delta$. Their joint projection is the scalar projection and vanishes on $\widehat X$. Therefore
\[
 1\le\sum_{j=1}^n\norm{(I-P_j)\widehat X}^2
 \le\frac{n(n+1)(2n+1)}6K_n^2\delta^2.
\]
This proves the bound. The two orientations $I-T_n^*T_n$ and $I-T_nT_n^*$ share the same singular defects. Stirling's formula gives the asymptotic exponent.
\end{proof}

\subsection{Consequence for arbitrary growing memory}
The witness has weight at most $m/2^n$ on the identity vector. Its energy and \eqref{eq:allgap} imply
\begin{equation}
 \frac m{2^n}\ge1-\frac{e_\epsilon}{g_n},\qquad
 q_n=\log_2m\ge n+\log_2(1-e_\epsilon/g_n),\label{eq:finiteq}
\end{equation}
when $e_\epsilon<g_n$. This bound is explicit jointly in $n,\theta,\epsilon$, but can be weak at moderate accuracy.

At fixed $\theta$, put $L_\epsilon=\log(1/\epsilon)$ and take
\[
 n=\left\lfloor\frac{(1-\eta)L_\epsilon}{4\log L_\epsilon}\right\rfloor,\qquad0<\eta<1.
\]
Then $g_n/\epsilon\to\infty$ and $q_n\ge n-o(1)$. Provided the stream reaches this cut,
\begin{equation}
 \liminf_{\epsilon\to0}\frac{q_{\rm peak}\log\log(1/\epsilon)}{\log(1/\epsilon)}\ge\frac14.\label{eq:asymemory}
\end{equation}
Logarithms here are natural; $q$ counts qubits. The branch-tail bounds give the same leading asymptotic for expected occupancy in the Bell-input experiment: use a threshold below $g_n$, where the low-energy mode count is one, then sum the memory tails.

At $N=\Theta_\theta(\log(1/\epsilon))$, including the quartic benchmark at fixed angle, this is $q_{\rm peak}=\Omega_\theta(N/\log N)$. It excludes $O(N^{1-\alpha})$ memory for every fixed $\alpha>0$ in that regime, with arbitrary classical memory and processing allowed. It still permits an $N/\log N$-type receiver. The angle dependence is not uniform: the explicit $g_n$ exposes it. No quartic memory lower bound as $\theta\to0$ at fixed error follows from this argument.
